% 第八章 总结与展望

\chapter{总结与展望}

\section{全文总结}

本文围绕FreeFlow项目，设计并实现了一个基于区块链与IPFS的链上音乐发行与播放系统。系统以桌面端应用为用户入口，以Web2.5后端作为业务服务和资源索引层，以Solidity智能合约作为链上访问控制和收益分账层，以IPFS/Pinata作为音乐资源存储层，形成了较完整的链上音乐业务闭环。

本文完成的主要工作包括：

\begin{enumerate}
  \item 设计了FreeFlow系统总体架构，明确桌面端、后端、合约和IPFS存储的职责分工。
  \item 实现了基于SIWE的钱包登录和Profile身份管理。
  \item 实现了面向创作者的音乐发布工作台，支持草稿管理、素材上传、metadata生成和链上发布。
  \item 实现了基于ERC-1155访问凭证的购买授权机制，并通过分账合约支持创作者收益分配。
  \item 实现了链上音乐资源检索与排序模块，支持文本字段和链上身份字段的统一检索。
  \item 实现了链上音乐详情、购买、播放、链上音乐库、评论和下载等用户侧功能。
\end{enumerate}

总体来看，FreeFlow更偏向系统设计与工程实现，链上音乐资源检索与排序算法作为系统中的辅助模块，用于增强资源发现能力和论文算法内容。

\section{不足与改进方向}

当前系统仍存在以下不足：

\begin{enumerate}
  \item 音频资源当前主要通过IPFS gateway访问，CID公开后仍可能绕过系统访问控制。后续可引入音频加密、授权网关或短期签名URL。
  \item 检索排序模块目前主要基于元数据、链上身份字段和聚合热度，缺少真实用户点击、播放和收藏反馈。后续可引入用户行为日志和个性化推荐。
  \item 后端当前主要保存链上发布结果投影，后续可增加链上事件监听器，自动同步发布、购买和分账事件。
  \item 系统测试仍需进一步完善，尤其是桌面端自动化测试、异常流程测试和多操作系统兼容性测试。
  \item 合约目前运行在测试网环境，若进入生产环境，还需要进行更严格的安全审计和gas成本优化。
\end{enumerate}

\section{后续展望}

后续可以从三个方向继续完善FreeFlow系统。第一，在安全方面引入加密音频和授权解密流程，使链上访问凭证真正控制完整资源访问。第二，在检索和推荐方面收集搜索、点击、播放、购买和评论行为，进一步提升排序质量。第三，在工程化方面完善链上事件同步、错误恢复、自动化测试和部署流程，使系统更加稳定可靠。

本章当前为初稿占位，后续应结合最终系统测试数据和论文整体内容进行压缩和精炼。
